\section{Prime patterns for the two conjugate local tests}
\label{pri:section}

This section proves Corollary~\ref{pri:integer-coverage}, the ordinary
coverage input used in the recursive theory.  The raw prime-pattern result
below supplies the reduction orders; its first consequence constructs the
paired integer witnesses by the Chinese Remainder Theorem.

Retain the five points $\mathcal C$ and the primitive integral forms
$L_b$ of Theorem~\ref{hei:height}, choosing $L_\infty(U,V)=V$.
For integers $A,u_1,u_2,v$ with $v\ne0$, setting $u_3=Av-u_1-u_2$ gives
\[
                         A=\frac{u_1}{v}+\frac{u_2}{v}+\frac{u_3}{v}.
\]
Evaluating the five contact forms on $(u_i,v)$ gives the common value $v$
and four further values for each slope. We seek to express each of these
thirteen entries as the product of a sign, a multiplier from one fixed
finite set, and a positive prime at which the reduction orders of $E$ and
$E^\sigma$ permit simultaneous local tests.

\begin{lemma}[Prime patterns]\label{pri:pattern}
There is a finite set $\mathcal K\subset\Z_{>0}$ such that, for every
$A\in\Z$ and every real $B>0$, there exist integers $u_1,u_2,v$, with
$v\ne0$, for which, on setting $u_3=Av-u_1-u_2$, each of the thirteen
numbers
\begin{equation}\label{pri:forms-list}
       v,\qquad L_b(u_j,v)
       \quad(1\le j\le3,\ b\in\mathcal C\setminus\{\infty\})
\end{equation}
has the form $\epsilon kr$, where $\epsilon\in\{1,-1\}$,
$k\in\mathcal K$, and $r>B$ is a positive rational prime.  Every such
prime splits in $F$ and is a prime of good reduction for both $E$ and
$E^\sigma$.  At either place of $F$ above $r$, if $d_1,d_2$ are the
cardinalities of the reduction groups of $E,E^\sigma$ over $\F_r$, then
\begin{equation}\label{pri:order-conditions}
                    \gcd(d_1,d_2)\mid4,
                    \qquad r\nmid d_1d_2.
\end{equation}
The set $\mathcal K$ is independent of $A$ and $B$.
\end{lemma}

The listed values are nonzero, so the slopes $(u_i:v)$ lie outside
$\mathcal C$. A prime counted by one of their contact radicals divides a
listed contact value, because the corresponding normalized contact depth
is positive and the coordinates are integral; unless it divides a member
of $\mathcal K$, it is therefore one of the selected primes $r$.

Before proving the lemma, we record the consequence that explains its
splitting and reduction-order conditions. Together with
Proposition~\ref{ell:formal}, they give two local tests at each place,
with one pair of auxiliary integers at that place serving every integer $a$.

\begin{corollary}[Simultaneous integer tests]\label{pri:integer-tests}
Fix an integer $K\ge1$ and a prime $r$ supplied by Lemma~\ref{pri:pattern}, large
enough for Proposition~\ref{ell:formal} at precision $K$.  At each of
the two places $v$ above $r$ there are nonzero integers $h,j_0$, chosen
independently of $a\in\Z$, such that, for every $a\in\Z$:
\begin{enumerate}
\item $hP$ and $haP$ reduce to $O$ at $v$,
      $\ell_K(z(hP))\ne0$, and
\[
           v\!\left(\frac{\ell_K(z(haP))}{\ell_K(z(hP))}-a\right)
                                                              \ge K;
\]
\item $j_0P^\sigma$ and $(h-4)aP^\sigma$ reduce to $O$ at the
      same place $v$, and
\[
            v\!\left(\frac{z((h-4)aP^\sigma)}{z(j_0P^\sigma)}\right)
                                                              \ge K.
\]
\end{enumerate}
The choices of $h,j_0$ may differ at the two places.
\end{corollary}

\begin{proof}
Let $d_1,d_2$ be the two reduction orders at the chosen place.  Choose
$j_0=d_2$.  Since $r\nmid d_1$ and $\gcd(d_1,d_2)\mid4$, the pair of
integer congruences
\[
                  h\equiv0\pmod{d_1},\qquad
                  h\equiv4\pmod{d_2r^K}
\]
is consistent.  Choose a positive solution $h$.  Reduction orders
show that all four indicated points are in the respective formal
kernels, uniformly for $a\in\Z$.  The logarithm congruence follows from
Proposition~\ref{ell:formal} with indices $ha,h$, whose ratio is $a$.
For the conjugate ratio use \eqref{ell:ratio} on the conjugate curve:
\[
 v\!\left(\frac{z((h-4)aP^\sigma)}{z(j_0P^\sigma)}\right)
          =v_r\!\left(\frac{(h-4)a}{d_2}\right)\ge K.
\]
The nonzero denominator follows because $j_0P^\sigma$ is nontorsion.
Zero numerator indices are covered by the infinite-valuation
convention.  Finally, at these split unramified primes the valuation
ring on $F$ is the localization of $\mathcal O_F=\Z[\sqrt2]$ at the
chosen prime ideal, and $r$ is a uniformizer.  Thus every displayed
inequality is exactly membership in $r^K(\mathcal O_F)_{\mathfrak p}$.
For a finite formal point, the condition $v(x)<0$ equivalently places
$1/x$ in $r(\mathcal O_F)_{\mathfrak p}$, as used in the ring-language tests of
Section~\ref{red:theory}.
\end{proof}

\begin{proof}[Proof of Corollary~\ref{pri:integer-coverage}]
Fix \(K\ge1\), \(A\in\Z\), and \(Y_0>0\).  Choose a real threshold
larger than \(Y_0\), \(K+1\), and every prime in the fixed finite set
\(S_E\).  Lemma~\ref{pri:pattern} gives the thirteen contact values
with primes above this threshold, using its one finite set
\(\mathcal K\).  These primes split in \(F\) and are good for both
curves.  At each of the two places above each occurrence,
Corollary~\ref{pri:integer-tests} gives nonzero \(h,j_0\), fixed for
every integer \(a\), with its two formal-kernel conditions and valuation
comparisons.  In the identity embedding with \(T_a=aP\), these are exactly
conditions (i) and (ii) of Lemma~\ref{red:local-comparison}.  The final
localization calculation in the proof of Corollary~\ref{pri:integer-tests}
gives the stated membership and finite-point interpretations.  The choices
are made separately at the different occurrences and places, as allowed.
The fixed set \(\mathcal K\) is independent of \(K,A,Y_0\) because the
prime-pattern lemma makes it independent of \(A\) and of the threshold.
\end{proof}

We now prove Lemma~\ref{pri:pattern}. Constants with a subscript $A$ may
depend on $A$, a subsequently fixed real box, and the fixed contact
forms, but not on the small-prime cutoff $L$. The order of choices is:
the finite multiplier set; then $A$, a box and one initial congruence
class; then $L$; and finally a large dilation parameter $X$.

\subsection{Reduction orders and their possible common factors}

Restrict attention to primes
\begin{equation}\label{pri:prime-residues}
                         r\equiv1\pmod8,
                         \qquad r\equiv3\pmod7.
\end{equation}
Quadratic reciprocity gives $(2/r)=1$, so $r$ splits in $F$, and
$(7/r)=(r/7)=-1$.  At either chosen place above $r$, the two nonzero
reductions of $d$ and $d^\sigma$ have product $7$.  They therefore
represent opposite square classes in $\F_r^*$.  The curves $E$ and
$E^\sigma$ are the quadratic twists by these classes of
$E_0:y^2=x^3-x$, so their traces of Frobenius are opposite.

Here is the specific description of that trace which we need.  Since
$r\equiv1\pmod4$, the automorphism
$(x,y)\mapsto(-x,iy)$ of $E_0$ is defined over $\F_r$ and has square
$[-1]$.  Its centralizer in the rational geometric endomorphism algebra
is $\Q(i)$.  Indeed, the faithful action of endomorphisms on an auxiliary
two-dimensional rational Tate module bounds this centralizer by
dimension two, while it already contains $\Q(i)$.  More explicitly one
may use the $3$-adic Tate module for $r>3$, on which the minimal
polynomial $T^2+1$ is irreducible and its matrix centralizer has dimension
two.  The faithful action remains injective after extending coefficients
\cite[Theorem III.7.4]{Silverman2009}.

Frobenius commutes with this automorphism.  Its characteristic polynomial
is integral, and its determinant is $r$; these are the usual degree and
trace identities for elliptic endomorphisms
\cite[Proposition III.8.6 and Chapter V, \S2]{Silverman2009}.
It is thus an algebraic integer in $\Q(i)$, hence equals $a+bi$ for
$a,b\in\Z$, with
\begin{equation}\label{pri:gaussian}
                    a^2+b^2=r,\qquad a\ne0.
\end{equation}
The last condition follows because a rational prime is not a square.
Consequently, in either order,
\begin{equation}\label{pri:orders}
                            d_1,d_2=r+1\pm2a.
\end{equation}
Both orders are divisible by four, since both curves have full rational
two-torsion over $\F_r$.  Their sum is $2(r+1)$, whose $2$-adic valuation
is exactly two by \eqref{pri:prime-residues}.  Thus the $2$-part of their
greatest common divisor is exactly four.  For sufficiently large $r$,
the Hasse bound gives $0<d_i<2r$.  Since $4\mid d_i$ and $r$ is odd,
$d_i\ne r$, proving $r\nmid d_i$.

If an odd prime $l$ divides both orders, their sum and difference show
that $l\mid r+1$ and $l\mid a$.  Equation \eqref{pri:gaussian} then gives
\begin{equation}\label{pri:bad-classes}
                    a\equiv0\pmod l,\qquad
                    b^2\equiv-1\pmod l.
\end{equation}
In particular such an $l$ is $1$ modulo $4$.  We shall first exclude
small $l$ by congruences on $r+1$, then use an upper sieve to remove the
remaining primes $r$ satisfying these two local conditions.

\subsection{A finite multiplier set and the initial residue conditions}

Write $L_b(U,V)=\alpha_bU+\gamma_bV$ for $b\ne\infty$.
Here $\alpha_b\ne0$, $\gcd(\alpha_b,\gamma_b)=1$, and
$\alpha_b\gamma_c-\alpha_c\gamma_b\ne0$ for $b\ne c$.
For an integer $A$, put $T=(u_1,u_2,v)$ and denote the thirteen
homogeneous forms in \eqref{pri:forms-list} by $F_\nu(T)$,
$1\le\nu\le t=13$. Their coefficient vectors, other than $(0,0,1)$,
are
\[
 (\alpha_b,0,\gamma_b),\quad
 (0,\alpha_b,\gamma_b),\quad
 (-\alpha_b,-\alpha_b,\alpha_bA+\gamma_b).
\]
Every vector is primitive, including the last one because
$\gcd(\alpha_b,\alpha_bA+\gamma_b)=1$. There is a fixed threshold
$p_0$, independent of $A$, above which any two vectors are linearly
independent over $\F_p$. Within one of the three families this follows
from the fixed nonzero determinants of the $L_b$; between different
families it follows from the three distinct directions
$(1,0),(0,1),(-1,-1)$ in the first two coordinates. The form $v$ is
independent of each remaining form whenever $p\nmid\alpha_b$.

Choose once and for all a finite set $S_0$ consisting of all primes up
to a sufficiently large threshold, including $2,7$, the exceptions
just described, and all fixed bad primes needed for the elliptic curves.
Increase the threshold so that $p>2t$ and the later local-density
estimates hold for every $p\notin S_0$. Choose a fixed integer $H_0$ with
$t2^{-H_0}<1$. For any $A$ and any $p\in S_0$, each primitive $F_\nu$
maps $(\Z/p^{H_0}\Z)^3$ onto $\Z/p^{H_0}\Z$. The proportion of vectors for
which $p^{H_0}\mid F_\nu(T)$ is therefore $p^{-H_0}$. Since $tp^{-H_0}<1$,
there is a vector class modulo $p^{H_0}$ on which none of the $F_\nu$ is
divisible by $p^{H_0}$. Every lift of this class modulo $p^{H_0+3}$ has
\[
                   v_p(F_\nu(T))<H_0\qquad(1\le\nu\le t).
\]
This proves the existence of the small-prime classes that we will
select for a particular $A$.

We next choose a finite catalogue of multipliers, without selecting
$A$ or any of those classes. Put $m_2=8$ and $m_p=p$ for odd
$p\in S_0$. Define units $c_2=1\pmod8$, $c_7=3\pmod7$, and
$c_p=1\pmod p$ for the remaining $p\in S_0$. These are the desired
residues of the positive primes in Lemma~\ref{pri:pattern}.
A catalogue datum consists
of integers and unit residues
\[
 e_{p\nu}\in\{0,\ldots,H_0-1\},\qquad
 \xi_{p\nu}\in(\Z/m_p\Z)^*
 \quad(p\in S_0,\ 1\le\nu\le t).
\]
For such a datum set $k_{0\nu}=\prod_{p\in S_0}p^{e_{p\nu}}$.
The Chinese Remainder Theorem and Dirichlet's Theorem
\cite[\S4]{Davenport2000} allow us to choose distinct primes
$h_1,\ldots,h_t\notin S_0$, hence above the fixed threshold, with
\begin{equation}\label{pri:multiplier-primes}
                 h_\nu\equiv c_p^{-1}\xi_{p\nu}\pmod{m_p}
                 \qquad(p\in S_0).
\end{equation}
When the datum later comes from a chosen vector class, $\xi_{p\nu}$ will
be the signed unit remaining after the factors from $S_0$ have been
removed. The congruence gives $\xi_{p\nu}h_\nu^{-1}\equiv c_p\pmod{m_p}$,
so division by the additional factor $h_\nu$ gives the desired residue.
There are only finitely many catalogue data. Fix one such tuple of
primes for each datum and let $\mathcal K$ be the union of the resulting
integers $k_{0\nu}h_\nu$. Thus $\mathcal K$ is finite and is fixed
before $A$, the cutoff $L$, and the dilation $X$.

Now fix the integer $A$ for which the lemma is to be proved. Choose a
bounded open rectangular box $\mathcal B$ of positive volume whose
closure avoids the zero planes of all $F_\nu$, and let
$\epsilon_\nu\in\{1,-1\}$ be the sign of $F_\nu$ on this box. For each
$p\in S_0$, choose one of the vector classes modulo $p^{H_0+3}$ supplied
above and put $e_{p\nu}=v_p(F_\nu(T))$ on that class. These values do
not depend on its representative. They define $k_{0\nu}$ as above.
The class also determines the unit
\[
 \xi_{p\nu}=
 \epsilon_\nu\frac{F_\nu(T)}{p^{e_{p\nu}}}
       \left(\frac{k_{0\nu}}{p^{e_{p\nu}}}\right)^{-1}
       \pmod{m_p}.
\]
Here the second factor is a unit modulo $m_p$, and the first quotient
is known modulo $p^{H_0+3-e_{p\nu}}$, a modulus containing at least four
powers of $p$. In particular it determines the required residue even
when $p=2$ and $m_p=8$. Select the tuple $h_\nu$ already fixed for this
datum, and put $k_\nu=k_{0\nu}h_\nu$.

At $p=h_\nu$, impose $F_\nu(T)=0\pmod p$. On this plane, each other
$F_j$ is nonconstant. Each condition
$\epsilon_jF_j(T)=0$ or $-k_j$ removes at most one line, and at most
$2(t-1)p<p^2$ points are removed. There is therefore a point satisfying
\[
          \epsilon_jF_j(T)\ne0,-k_j\pmod{h_\nu}
          \quad(j\ne\nu).
\]
Lift that point modulo $p^2$ so that
\[
                    \epsilon_\nu F_\nu(T)/p
                             \ne0,-k_{0\nu}\pmod p.
\]
Such a lift exists because $F_\nu$ is primitive, so varying the lift
changes $F_\nu(T)/p$ through all residues modulo $p$. The Chinese
Remainder Theorem combines these choices with the chosen classes at
$S_0$ into one vector class modulo
\begin{equation}\label{pri:initial-modulus}
              W_0=\prod_{p\in S_0}p^{H_0+3}
                          \prod_{\nu=1}^t h_\nu^2.
\end{equation}
On this class all
\[
                          \Psi_\nu(T)=\epsilon_\nu F_\nu(T)/k_\nu
\]
are integers and are nonzero modulo every prime dividing $W_0$.
At a prime in $S_0$, the definition of $\xi_{p\nu}$ and
\eqref{pri:multiplier-primes} give $\Psi_\nu\equiv c_p\pmod{m_p}$.
The conditions at each $h_j$ give the remaining exclusions. Thus the
divided forms equal $1$ modulo $8$ and $3$ modulo $7$, and avoid $-1$
modulo every odd prime dividing $W_0$. Notice also that $W_0/k_\nu$ is
divisible by every prime dividing $W_0$. This extra divisibility will
control the local factors of the prime-counting theorem.

Fix a smaller rectangular box $\mathcal B'$ of positive volume with
$\overline{\mathcal B'}\subset\mathcal B$. Since the signs agree with
the forms on $\mathcal B$ and its closure avoids their zero planes,
there are constants $0<c_A<C_A$ such that
\begin{equation}\label{pri:value-range}
        c_AX\le\Psi_\nu(T)\le C_AX
          \quad(T\in X\mathcal B,\ 1\le\nu\le t).
\end{equation}
All choices for this $A$, including $W_0$, are now fixed before choosing
$L$. The catalogue was finite, so $W_0$ ranges in a fixed finite set as
$A$ varies.

\subsection{A lower bound for the prime patterns}

Let $L$ be a fixed cutoff larger than every prime dividing $W_0$, and put
\[
                    W=W_0\prod_{\substack{p\le L\\p\nmid W_0}}p.
\]
At each new prime $p\le L$ allow all residue classes of $T$ satisfying
\begin{equation}\label{pri:small-exclusions}
                   \epsilon_\nu F_\nu(T)\ne0,-k_\nu\pmod p
                         \quad(1\le\nu\le t).
\end{equation}
At most $2t$ affine planes, each of $p^2$ points, are excluded.  Thus if
$\mathcal R_W$ is the set of all allowed vector classes modulo $W$,
including the one fixed class modulo $W_0$, then
\begin{equation}\label{pri:crt-density}
 \frac{|\mathcal R_W|}{W^3}
 \ge W_0^{-3}\prod_{\substack{p\le L\\p\nmid W_0}}(1-2t/p)
 \ge c_0(\log L)^{-C_0}.
\end{equation}
The last bound follows by taking logarithms and using Mertens' estimate
$\sum_{p\le L}p^{-1}=\log\log L+O(1)$
\cite[equation~(1.1) and Theorem~1]{Williams1974}.  The constants can be fixed independently of
$L$; since $W_0$ ranges in a fixed finite set they could also be made
independent of $A$.  Whenever $\Psi_\nu(T)=r$ is prime, these conditions
and the initial conditions imply that no odd $l\le L$ divides $r+1$.
Such an $l$ cannot divide both orders in \eqref{pri:orders}.

Let $N_{A,L}(X)$ be the number of $T\in X\mathcal B\cap\Z^3$ for which
every $\Psi_\nu(T)$ is prime and all imposed congruences hold. We claim
that
\begin{equation}\label{pri:lower}
                  N_{A,L}(X)\ge
        c_A(\log L)^{-C_0}\frac{X^3}{(\log X)^t}
\end{equation}
for sufficiently large $X$ depending on $A,L$, with $c_A>0$ independent
of $L$.

We use the following precise form of the Green--Tao--Ziegler theorem.
For a fixed finite collection of integer affine-linear forms
$\psi_1,\ldots,\psi_t$ on $\Z^d$ whose nonzero linear parts are pairwise
nonproportional, and a fixed dilating box where all values are positive,
the von Mangoldt weighted count is
\begin{equation}\label{pri:gtz}
 \sum_{n\in\mathcal D_X\cap\Z^d}\prod_{\nu=1}^t\Lambda(\psi_\nu(n))
 =\operatorname{vol}(\mathcal D_X)\prod_p\beta_p+o(X^d),
\end{equation}
where
\[
 \beta_p=(1-1/p)^{-t}p^{-d}
   \#\{n\in\F_p^d:\psi_\nu(n)\ne0\text{ for every }\nu\}.
\]
Here all coefficients and the box are fixed before $X\to\infty$.
This is the finite-complexity case of the Main Theorem of
\cite{GreenTao2010}, combined with the M\"obius--nilsequence theorem
\cite[Theorem 1.1]{GreenTao2012} and the inverse theorem
\cite[Theorem 1.3]{GreenTaoZiegler2012}, with its correction
\cite{GreenTaoZiegler2024Erratum}.  Pairwise nonproportional linear parts
give finite complexity: for any chosen form, the other forms can be
partitioned into singletons, none spanning it.  The 2024 erratum
corrects intermediate factorization and filtration statements and leaves
the stated inverse theorem, and hence \eqref{pri:gtz}, unchanged.

For each $T_0\in\mathcal R_W$, use its representative in $[0,W)^3$ and
write $T=T_0+Wn$.  The forms
\[
             \psi_{\nu,T_0}(n)
              =\epsilon_\nu F_\nu(T_0+Wn)/k_\nu
\]
have integer coefficients and pairwise nonproportional linear parts.
Use $n\in(X/W)\mathcal B'$.  For all sufficiently large $X$, depending
on $A,L$, this entails $T\in X\mathcal B$, uniformly over the finite
set of representatives $T_0$.  The volume in \eqref{pri:gtz} is then
$\operatorname{vol}(\mathcal B')X^3/W^3$.

At $p\mid W$, the coefficients of the variable parts of every divided
form vanish modulo $p$, whereas their constant values are nonzero.
Consequently $\beta_p=(1-1/p)^{-t}\ge1$.  At $p\nmid W$, multiplication
by $W$ and by $k_\nu^{-1}$ is invertible.  The zero sets of the forms
are distinct affine hyperplanes, and each pair has an intersection of
codimension two.  Inclusion--exclusion, using only the first two terms
for a uniform error, gives
\[
        p^{-3}\#\{n:\psi_{\nu,T_0}(n)\ne0\ (\forall\nu)\}
                            =1-t/p+O_t(p^{-2}).
\]
Thus $\beta_p=1+O_t(p^{-2})$.  Taking the original fixed threshold
large enough, the product over any subset of these primes is bounded
below by one fixed positive convergent product.  This lower bound is
independent of $L$ and $T_0$.

Prime powers in \eqref{pri:gtz} are negligible.  On a box of scale $X$,
there are $O(X^{1/2}\log X)$ proper prime powers in the range of each
nonconstant form, and at most $O_{A,L}(X^2)$ lattice points for each
prescribed value of that form, by solving for one coordinate.  Their
total weighted contribution is therefore
$O_{A,L}(X^{5/2}(\log X)^{t+1})=o_{A,L}(X^3)$.
For genuine prime values the weight is at most $C_A(\log X)^t$, by
\eqref{pri:value-range}.

Sum \eqref{pri:gtz} over the finitely many allowed progressions. For
each fixed $L$ their errors still sum to $o_{A,L}(X^3)$, while the
$W^{-3}$ volume factor cancels the $W^3$ in the number of progressions
in \eqref{pri:crt-density}. The lower bound for the local-factor product,
the removal of prime powers, and the bound on the prime weights therefore
give \eqref{pri:lower}, with a positive coefficient independent of $L$.
We require no uniformity in the little-oh error as $A$ or $L$ varies.

We must now remove the patterns for which some pair of reduction orders
has an odd common factor $l>L$. For one form, we will bound the number of
possible bad prime values by $O_A(X/(L\log X))$ plus terms of smaller
order as $X\to\infty$ for fixed $L$. For each fixed bad value, a second
count will bound the points on
its affine fiber for which the other $t-1$ forms remain prime by
$O_A(X^2/(\log X)^{t-1})$. The contribution from the term
$X/(L\log X)$ is therefore $O_A(X^3/(L(\log X)^t))$. This retains the
factor $(\log X)^{-t}$ in \eqref{pri:lower} and introduces the saving $1/L$.
The first count sieves the pairs $(a,b)$ in \eqref{pri:gaussian}; the
second sieves the two free coordinates on a fiber. Both use the following
two-dimensional estimate.

\subsection{An elementary upper sieve with uniform errors}

Let $\mathcal P_Z$ be a set of primes at most $Z$, with forbidden subsets
$\Omega_p\subset\F_p^2$ having densities
$g(p)=|\Omega_p|/p^2$ in $(0,1/2)$.  Extend $g$ multiplicatively to
squarefree integers supported on $\mathcal P_Z$. Let $\mathbf x$ range
over a square in $\Z^2$ with $n$ integer choices in each coordinate.
The simultaneous forbidden conditions at the primes dividing a
squarefree integer $e$ hold at
\begin{equation}\label{pri:residue-error}
                          n^2g(e)+O(ne+e^2)
\end{equation}
points, uniformly in the location of the square.  Indeed the Chinese
Remainder Theorem gives $e^2g(e)$ residue classes, and the count in
each differs from $n^2/e^2$ by $O(n/e+1)$.

The following estimate uses Selberg's $\lambda^2$ method
\cite{Selberg1947}; compare \cite[Section~3.2]{MontgomeryVaughan2007}.
We give the weights and error term needed for these two-dimensional boxes.

\begin{lemma}\label{pri:upper-sieve}
The number of points in this square avoiding every $\Omega_p$ is at most
\begin{equation}\label{pri:sieve-bound}
            \frac{n^2}{G(Z)}+O(nZ^6+Z^8),\qquad
 G(Z)=\sum_{\substack{e\le Z\text{ squarefree}\\p\mid e\Rightarrow p\in\mathcal P_Z}}
                           \prod_{p\mid e}\frac{g(p)}{1-g(p)}.
\end{equation}
The error constant is absolute.
\end{lemma}

\begin{proof}
All integers in the sums that follow are squarefree and supported on
$\mathcal P_Z$.  Put $g_*(p)=(1-g(p))/g(p)$ and extend it
multiplicatively, with $g(1)=g_*(1)=1$.  For real weights $\lambda_e$
supported on $e\le Z$ with $\lambda_1=1$, the square
\[
       \left(\sum_{\substack{e\le Z\\
                 \mathbf x\bmod p\in\Omega_p\ (p\mid e)}}\lambda_e\right)^2
\]
majorizes the indicator of avoiding all forbidden sets: at such a point
only $e=1$ contributes.  Upon summing and applying
\eqref{pri:residue-error}, its main term divided by $n^2$ is
$\sum_{e,f}\lambda_e\lambda_fg([e,f])$.
The identity $1/g(d)=\sum_{h\mid d}g_*(h)$ gives the diagonalization
\begin{equation}\label{pri:diagonalization}
 \sum_{e,f}\lambda_e\lambda_fg([e,f])
   =\sum_hg_*(h)
            \left(\sum_{\substack{e\le Z\\h\mid e}}\lambda_eg(e)\right)^2.
\end{equation}
Prescribe the inner sums to equal
$\mu(h)/(G(Z)g_*(h))$.  Finite M\"obius inversion gives
\begin{equation}\label{pri:weights}
 \lambda_e=\frac{\mu(e)}{G(Z)g(e)g_*(e)}
       \sum_{\substack{k\le Z/e\\(k,e)=1}}
                       \frac{\mu(k)^2}{g_*(k)}.
\end{equation}
In particular $\lambda_1=1$ and
\[
             |\lambda_e|\le\frac1{g(e)g_*(e)}
                  =\prod_{p\mid e}(1-g(p))^{-1}
                  \le2^{\omega(e)}\le e.
\]
Substitution in \eqref{pri:diagonalization} gives $1/G(Z)$.
For the error at $[e,f]$, use $[e,f]\le ef$ and
$|\lambda_e\lambda_f|\le ef$.  Its sum is bounded by
\[
     O\!\left(n\sum_{e,f\le Z}e^2f^2+
                       \sum_{e,f\le Z}e^3f^3\right)
                           =O(nZ^6+Z^8),
\]
as required.
\end{proof}

We shall always set $Z=X^{1/100}$.  Two consequences for $G(Z)$ will
be useful.  If $\mathcal P_Z$ contains all primes above a fixed
threshold up to $Z$, and
$g(p)=s/p+O_s(p^{-2})$ for a fixed positive integer $s$, then
\begin{equation}\label{pri:sieve-dimension}
                              G(Z)\gg(\log X)^s.
\end{equation}
If instead it contains the primes $p\equiv1\pmod4$ above a fixed
threshold up to $Z$, with at most one additional prime omitted, and
$g(p)=2/p+O(p^{-2})$, then
\begin{equation}\label{pri:split-sieve-dimension}
                              G(Z)\gg\log X,
\end{equation}
uniformly in that omitted prime.  In these two estimates the implicit
constants may depend on the fixed threshold and the fixed density
estimates, as well as on $s$, but not on $X$ or the omitted prime.

Here are details of the truncation behind both claims.  For
$y=Z^\theta$, where $\theta>0$ will be fixed sufficiently small, the
Euler product
\[
             D(y)=\prod_{\substack{p\le y\\p\in\mathcal P_Z}}
                               (1+g_*(p)^{-1})
\]
is respectively $\gg(\log y)^s$ or $\gg\log y$, by the ordinary and
progression Mertens estimates
\cite[equation~(1.1) and Theorem~1]{Williams1974}.  In the latter case removing one prime
loses at most a factor two because $g(p)<1/2$.
In the expansion of $D(y)$ give a squarefree $e$ weight $1/g_*(e)$.
The weighted mean of $\log e$ equals
\[
                          \sum_{\substack{p\le y\\p\in\mathcal P_Z}}
                                 g(p)\log p\ll\log y.
\]
This uses $\sum_{p\le y}(\log p)/p\ll\log y$, and the bound is uniform
when a prime is omitted.  For a small enough fixed $\theta$, Markov's
inequality shows that at least half the weight has $e\le Z$.
Those terms occur in $G(Z)$, proving both assertions.

\subsection{Counting bad prime values}

Fix one index $\nu$.  All its possible prime values lie in the interval
\eqref{pri:value-range}.  Let $B_{A,L}(X)$ be the number of primes in
that interval satisfying \eqref{pri:prime-residues} for which the two
orders have a common odd factor $l>L$.  We claim
\begin{equation}\label{pri:bad-primes}
            B_{A,L}(X)\ll_A
             \frac{X}{L\log X}+X^{7/8}+\sqrt X\log X.
\end{equation}
It suffices to count pairs $(a,b)\in\Z^2$ with
$|a|,|b|\le C_A\sqrt X$, $a\ne0$, for which $a^2+b^2$ is a prime
in the interval and \eqref{pri:bad-classes} holds for some prime
$l>L$.  Every bad prime supplies such a pair.  Overcounting pairs, and
counting a pair several times if necessary, gives an upper bound.

First consider $L<l\le X^{1/8}$.  If $-1$ has no square root modulo
$l$, there are no pairs.  Otherwise take each of its two roots $b_0$
and write
\[
                          a=lU,\qquad b=b_0+lV.
\]
The relevant $U,V$ lie in an encompassing square with
\begin{equation}\label{pri:short-square}
               n\asymp_A\sqrt X/l\gg_A X^{3/8}
\end{equation}
integer choices in each coordinate, uniformly for these $l$ and for
least nonnegative representatives $b_0$.
Sieve by all sufficiently large primes $p\equiv1\pmod4$, $p\ne l$,
up to $Z=X^{1/100}$.  Modulo $p$ the substitution
$(U,V)\mapsto(lU,b_0+lV)$ is an affine bijection, so the forbidden
equation $a^2+b^2=0$ has exactly $2p-1$ solutions: it is the union of
two distinct lines meeting at one point.  Thus
\[
                            g(p)=2/p-1/p^2,
\]
independently of $l,b_0$.  Since the prime value $a^2+b^2$ is
$\asymp_A X>Z$, it avoids all these forbidden sets.
Lemma~\ref{pri:upper-sieve} and
\eqref{pri:split-sieve-dimension} give
$O(n^2/\log X+nZ^6+Z^8)$ pairs.  At the shortest possible squares,
\[
 \frac{nZ^6+Z^8}{n^2/\log X}
       \ll_A(\log X)
                   \bigl(X^{6/100-3/8}+X^{8/100-3/4}\bigr)=o_A(1).
\]
The error is therefore absorbed uniformly before summing over $l$.
Each $l$ contributes $O_A(X/(l^2\log X))$, and summing over even all
integers $l>L$ gives $O_A(X/(L\log X))$.

For $l>X^{1/8}$ use a lattice count without a sieve.  Since $a\ne0$
and $l\mid a$, necessarily $l\le C_A\sqrt X$, and the number of
possible $a$ is $O_A(\sqrt X/l)$, with no additive constant.
There are at most two classes for $b$, giving
$O_A(\sqrt X/l+1)$ possibilities.  The count for one $l$ is thus
\[
                         O_A(X/l^2+\sqrt X/l).
\]
To sum it explicitly, split $X^{1/8}<l\le C_A\sqrt X$ into dyadic
blocks $D<l\le2D$.  Even if every integer in a block were allowed, its
contribution would be $O_A(X/D+\sqrt X)$.  The first terms form a
geometric sum $O_A(X^{7/8})$, while there are $O_A(\log X)$ blocks.
Their total is $O_A(X^{7/8}+\sqrt X\log X)$.  This proves
\eqref{pri:bad-primes} with constants independent of $L$.

\subsection{A uniform sieve on each affine fiber}

For a fixed prime value $r\asymp_A X$ of the $\nu$th form, consider
\begin{equation}\label{pri:fiber}
                         F_\nu(T)=\epsilon_\nu k_\nu r.
\end{equation}
We claim that the number of $T\in X\mathcal B\cap\Z^3$ on this
fiber for which all other $\Psi_j(T)$ are primes is
\begin{equation}\label{pri:fiber-bound}
                              \ll_A\frac{X^2}{(\log X)^{t-1}},
\end{equation}
uniformly in $r$ and independently of $L$.  In this upper bound we
discard all the congruence conditions used for the lower bound.

Since $F_\nu$ is primitive, an integral unimodular coordinate change
makes it the first coordinate.  The other two coordinates of every
point of the fiber inside $X\mathcal B$ lie in an encompassing square
with $n\asymp_A X$ integer choices in each direction.  The coordinate
change depends on $A,\nu$ but not on $r,L$.
Sieve by all primes up to $Z=X^{1/100}$ above a fixed sufficiently
large threshold, excluding every prime dividing any multiplier.
For all sufficiently large $X$, every sieving prime differs from $r$,
so the fiber constant
$c=\epsilon_\nu k_\nu r$ is nonzero modulo each such prime $p$.

On the affine plane $F_\nu=c$ over $\F_p$, every other $F_j$ restricts
to a nonconstant affine function, because $F_j$ is not proportional to
$F_\nu$.  Furthermore the $t-1$ zero lines of these restrictions are
distinct.  If two coincided, their affine functions would be scalar
multiples on that plane.  For some $\lambda\in\F_p^*$ and
$\mu\in\F_p$ this would give the identity
\[
                         F_j-\lambda F_k=\mu(F_\nu-c).
\]
The left side is homogeneous, so comparison of constant terms gives
$\mu c=0$.  Since $c\ne0$, one has $\mu=0$, contradicting the pairwise
independence of $F_j,F_k$.  This argument also covers the case that
three ambient forms are jointly dependent: a nonzero constant shift
then gives distinct parallel lines, rather than coincident ones.

The union of the $t-1$ distinct zero lines has
$(t-1)p+O_t(1)$ points.  All must be avoided, since the other prime
values have size $\asymp_A X$ and exceed the sieving primes; the
multiplier is invertible modulo every sieving prime.  Thus
$g(p)=(t-1)/p+O_t(p^{-2})$, uniformly in the fiber.  Lemma
\ref{pri:upper-sieve} and \eqref{pri:sieve-dimension} give
\[
       O_A\!\left(\frac{X^2}{(\log X)^{t-1}}+XZ^6+Z^8\right)
                       =O_A\!\left(\frac{X^2}{(\log X)^{t-1}}\right),
\]
which proves \eqref{pri:fiber-bound}.  All excluded sieving primes were
fixed in advance of $r,L$, and residue-count errors are uniform in the
position of the encompassing square.  This proves the claimed
uniformity.

\subsection{Completing the prime-pattern construction}

For each of the $t$ forms, multiply the number of its possible bad prime
values in \eqref{pri:bad-primes} by the uniform fiber bound
\eqref{pri:fiber-bound}.  A union bound shows that at most
\begin{equation}\label{pri:total-loss}
          \frac{C_AX^3}{L(\log X)^t}
                         +o_A\!\left(\frac{X^3}{(\log X)^t}\right)
\end{equation}
of the prime patterns can have a common odd factor $l>L$ in one of
their pairs of reduction orders.  More explicitly, the ratios of the
two other terms to $X^3/(\log X)^t$ are
$O_A(X^{-1/8}\log X)$ and $O_A(X^{-1/2}(\log X)^2)$; both tend to
zero.  None of the constants in this upper estimate depends on $L$.

Choose $L$, after all data depending on $A$ have been fixed, large
enough that
\[
                    C_A/L<\tfrac12 c_A(\log L)^{-C_0}.
\]
This is possible because $(\log L)^{C_0}/L\to0$.  Then choose $X$
sufficiently large for \eqref{pri:lower}, for all the upper estimates,
and for the residual little-oh in \eqref{pri:total-loss} to be smaller
than the remaining lower-bound coefficient.  At least one prime
pattern survives.  Small common odd factors were already excluded by
the CRT conditions, and all larger ones have just been removed.  The
$2$-part and the prime-to-$r$ assertions proved from
\eqref{pri:orders} now give \eqref{pri:order-conditions}.  The primes
have size at least $c_AX$, so increasing $X$ makes all of them exceed
any prescribed $B$.  This proves Lemma~\ref{pri:pattern}.

Together with Corollary~\ref{pri:integer-tests}, the prime-pattern lemma
proves Corollary~\ref{pri:integer-coverage}, completing the ordinary integer
witnesses used in Section~\ref{sec:reduction}.  The analysis there combines
the tested representations with the parity, height, and elliptic inputs to
turn success of all finite tests into an ordinary integer zero.
